\section{Hadronic tau leptons}
\label{sec:ch6_tau}

\subsection{Reconstruction}

$\tau$ leptons can decay both leptonically and hadronically.
Approximately 65\% of $\tau$ decays are hadronic, with the major modes containing charged and possible nerutral pions~\cite{ATLAS2022TauR22}.
In this anlysis we only consider the hardronic decay tau. It is denoted by $\tau_{\mathrm{had}}$. Here the $\tau$-neutrino is exculded in the reconstruction.

Hadronic tau candidates are first seeded by anti-$k_t$ jets with $R=0.4$.
To identify the $pp$ collision associated with each candidate, tracks within $\Delta R<0.2$ of the seed axis are then used to select the production vertex, choosing the vertex with the largest fraction of their $\pT$ sum.
This vertex provides the reference for track association and impact-parameter calculations.
The associated tracks are subsequently classified by a recurrent neural network (RNN) to select tau-decay tracks.
The number of selected tau-decay tracks defines the prong multiplicity: candidates with one or three such tracks are called one-prong or three-prong candidates, respectively.

\subsection{Identification}

Hadronic tau decays are identified using an RNN discriminant to distinguish them from quark- and gluon-initiated jets~\cite{ATLAS2019TauRNN,ATLAS2022TauR22}.
Track momenta and calorimeter shower properties are used as inputs, since tau decays typically have narrower showers and lower charged multiplicity than these jets.
Track displacements and, for three-prong candidates, the significance of the reconstructed decay-vertex displacement are also included as inputs sensitive to the finite tau lifetime.
Nearby activity is also included in the RNN inputs, so isolation information is used as part of the identification.

Baseline tau candidates are required to pass \texttt{Loose} RNN identification, and signal candidates must additionally satisfy \texttt{Tight}.
The higher RNN-score threshold for \texttt{Tight} rejects more jets at the cost of tau efficiency.
For reconstructed one-prong (three-prong) candidates, the target identification efficiencies are 85\% (75\%) for \texttt{Loose} and 60\% (45\%) for \texttt{Tight}~\cite{ATLAS2022TauR22}.

Electrons can be misidentified as hadronic tau candidates because their narrow showers and associated tracks resemble the signature of a one-prong tau decay~\cite{ATLAS2022TauR22}.
To suppress this background, both baseline and signal tau candidates are required to pass \texttt{Loose} electron rejection.
This working point requires a minimum score from a separate RNN using shower and track information, including the relation between calorimeter energy and track momentum.
The tau efficiencies are approximately 95\% and 98\% for one-prong and three-prong candidates, respectively.
The corresponding electron rejection efficiencies are approximately 99.7\% and 98.8\%~\cite{ATLAS2022TauR22}.

\subsection{Calibration}

The tau energy scale (TES) calibration is derived from MC samples~\cite{ATLAS2022TauR22}.
In these samples, the generated visible tau $\pT$, excluding the decay neutrino, is used as the reference for correcting the reconstructed $\pT$.
The reconstructed $\pT$ is measured both from calorimeter deposits alone and from a combination of charged-hadron track momenta and neutral-particle calorimeter deposits.
These two measurements are combined in a weighted average.
The averaged $\pT$ is then corrected using a boosted regression tree trained on these MC samples.

In addition to the momentum calibration, differences in tau reconstruction, identification and electron-rejection efficiencies between data and simulation must be correcte.
For this purpose, data-to-simulation scale factors are applied to simulated event weights.
In this analysis, for signal candidates passing \texttt{Tight} tau identification, the electron-veto scale factors associated with \texttt{Medium} tau identification are used because no corresponding factors were available for \texttt{Tight}.
